\subsubsection{Why Map QUBO to Quantum Mechanics?}\label{sec:why-map-qubo-to-quantum-mechanics}

At this point we have two seemingly unrelated descriptions of a problem.

\highspace
On one side, we have a \textbf{QUBO optimization problem} (\autopageref{sec:introduction-to-qubo}):
\begin{equation*}
    \argmin_{\mathbf{x}} E(\mathbf{x}) = \argmin_{\mathbf{x}} \mathbf{x} Q \mathbf{x}^T
\end{equation*}
Where the \hl{goal is to find the binary assignment $\mathbf{x}$ that minimizes the energy function $E(\mathbf{x})$}. Here, we are still in the classical mathematical formulation of the problem, and we have not yet made any connection to quantum mechanics, QUBO is not a quantum concept, but rather a classical optimization problem.

\highspace
On the other side, quantum mechanics describes physical systems through a Hamiltonian $\hat{H}$ (\autopageref{sec:hamiltonian-recap}), whose \hl{eigenvalues represent the possible energy levels of the system}.

\highspace
The main observation is that both descriptions revolve around the same concept: \textbf{energy}.

\highspace
\textcolor{Green3}{\faIcon[regular]{lightbulb} \textbf{The Central Idea.}} Recall what we do when solving a \textbf{QUBO problem}. We assign an energy value to every possible configuration (e.g., state $00$ has energy $E(00)$, state $01$ has energy $E(01)$, etc.). The \textbf{solution} is simply the state with the \textbf{lowest energy}. Now, consider a quantum system. It also possesses energy levels $E_0, E_1, E_2, \ldots$, and naturally tends to occupy its lowest-energy configuration, called the \textbf{ground state} (\autopageref{def:ground-state}), because it is the most stable. Therefore, we can think of a quantum system as an \textbf{energy minimization process}, where the system evolves towards its ground state, which corresponds to the minimum energy configuration.
\begin{equation*}
    \text{Quantum System} = \text{Energy Minimization Process}
\end{equation*}

\highspace
\textcolor{Green3}{\faIcon{check-circle} \textbf{The Opportunity.}} This similarity (both QUBO and quantum systems are about energy minimization) opens up a powerful opportunity: \textbf{we can leverage the natural energy-minimizing behavior of quantum systems to solve QUBO problems}. Precisely, instead of solving the QUBO problem directly using a classical algorithm, we can:
\begin{enumerate}
    \item Encode the QUBO energy landscape into a Hamiltonian.
    \item Build a quantum system described by that Hamiltonian.
    \item Let the quantum system search for its ground state.
    \item Read the solution from the resulting quantum state.
\end{enumerate}
Therefore, instead of viewing the optimization problem as a purely mathematical object, we reinterpret it as a physical system. The optimization then becomes:
\begin{equation*}
    \text{Find minimum of } E(\mathbf{x}) \quad \Rightarrow \quad \text{Find ground state of } \hat{H}
\end{equation*}
Or more precisely:
\begin{equation*}
    \text{Minimizing a QUBO} \iff \text{Finding the ground state of a Hamiltonian}
\end{equation*}